\cotefile{157-5}
\scanpage{648}\mspage{36}
les foncteurs $u^{W}_{!}$. Notons déjà que \unsure{la théorie|le théorème} de
K.S. Brown, \inserted{sur les \unsure{instruments} de subdivision\textsuperscript{$*$},} donne essentiellement ceci :

% slanted
{\reversemarginpar\margin{\unsure{revus|revu} \deleted{\unsure{et}} \unsure{généralisés|généralisé} par mes soins}}

\uline{Proposition 3.2}. \deleted{Le prédérivateur défini par
$\mathrm{Ho}_{W}$, quand} Soit $W\inserted{\subset\mathrm{Fl}\,M}$ \deleted{est associé} un localiseur
associé à une structure à cofibrations
$(C,W)$ sur $M$. \deleted{\illegible{1}} \textcircled{\scriptsize 1$^\circ$} Sur la catégorie $\mathrm{Diag}_0$ des
ensembles ordonnés \uline{finis}, ce prédérivateur
satisfait aux axiomes Der 1 (décomposition),
Der 2 (localisation), Der 3 (objet final),
Der 4 bis (existence de $u^{W}_{!}$), Der 5 bis
(calcul des fibres de $u^{W}_{!}$), Der 6 bis
(ex. \unsure{strictes|suites} à gauche). \textcircled{\scriptsize 2$^\circ$} C'est encore vrai
sur la \inserted{\unsure{2-}}catégorie $\mathrm{Diag}_1$ \inserted{\unsure{de} Cat plus grande,} formée des ens.
ordonnés $I$ \deleted{\uline{noethériens}} tels que l'espace
top. associé à $I^\circ$ soit noethérien (\unsure{i.e.} tte
suite croissante d'objets de $I$, ou d'ouverts
de $I^\circ\deleted{\unsure{\}}}$, est stationnaire). \textcircled{\scriptsize 3$^\circ$} Enfin, si
\deleted{\unsure{$I$}} l'axiome transfini \inserted{Cof 4} est satisfait i.e.
si $C$ et $TC$ sont stables par composition
transfinie, alors les conclusions de $1^\circ$ restent
valables sur la \deleted{\unsure{cat.}} sous-catégorie \inserted{pleine} $\mathrm{Diag}_2$ de
Cat formée des ens ordonnés qui sont
noethériens sans plus.

\uline{Corollaire 3.3} Supposons que l'on ait Cof 3 sous
la forme plus forte d'une factorisation
% archive sheet number: 647
